\section*{NS91: combine before estimating}

\paragraph{Status and classification.}
Exact identities, proved complete-tail inequalities, and finite certified computations.
This continues PR42 and PR45 and integrates the previously standalone draft.
No cofinal arithmetic lower gain is proved.

Retain the notation of NS90:
\[
 R=\tau-\Pi_N\tau,\quad E=\|R\|^2,\quad
 \alpha=\log2,\quad D=\alpha\tau-\sum_{n\le2N}d_na_n,
 \quad \ell=\alpha E-\langle R,D\rangle.
\]
Here $d_n=-\mu(n)\alpha$ through $N$, and
$d_n=-\mu(n)\log(2N/n)$ on the next block. Its complete projected cost
$K=\|(I-\Pi_N)\sum d_na_n\|^2$ is the NS73 cost; it is not replaced
by a model cost. Set $s=-\sum d_n/n$.

\paragraph{Exact first-primitive identity.}
Define
\[
 Q(t)=\int_t^\infty R(u)\frac{du}{u^2},\qquad J(t)=tD'(t).
\]
For $t\ge1$, away from immaterial endpoint values,
\[
 J(t)=\alpha-\sum_{n\le2N}d_n\{t/n\}
     =st+\sum_{k\le t}\delta_k,
 \qquad \delta_k=\alpha\mathbf1_{k=1}+\sum_{n\mid k}d_n.
\]
For $0<t<1$, $J(t)=st$.
Thus $J$ is bounded for $t\ge1$, even though the separate terms of the
NS90 divisor-potential sum have divergent absolute allowance. Since
$Q'=-R/t^2$, integration by parts gives the absolutely convergent identity
\begin{equation}\label{ns91-first}
 \langle R,D\rangle=\int_0^\infty J(t)Q(t)\frac{dt}{t}.
\end{equation}
The endpoints vanish: $D=O(t)$ and $Q=O(1+|\log t|)$ at zero, whereas
$D=O(1+\log t)$ and $Q=O((1+\log t)/t)$ at infinity. These bounds also
prove absolute integrability after this regrouping. No limit of separately
divergent terms is used.

\paragraph{Exact old-space correction: remove the whole initial interval.}
NS90 gives $D(t)=st$ for $0<t\le N$. Set
\[
 \widetilde D=D-Ns\,a_N,\qquad
 \widetilde d_N=d_N+Ns=-N\sum_{n\ne N}d_n/n,
 \qquad \widetilde d_n=d_n\ (n\ne N).
\]
Then $\sum\widetilde d_n/n=0$, and $a_N(t)=t/N$ on $0<t\le N$, so
\begin{equation}\label{ns91-zero-initial}
 \widetilde D(t)=\widetilde J(t)=0\quad(0<t<N).
\end{equation}
Because $a_N\in V_N$, $\langle R,\widetilde D\rangle=\langle R,D\rangle$.
The projected direction and its cost $K$ are also unchanged. This is a
change of the representative used to estimate the numerator, not a change
of the underlying gain direction or a refit. The subsequent formulas
apply to either coefficient vector; the corrected one is used for the
stronger certificates.

\paragraph{Complete dyadic tail inequality.}
Suppose $T\ge2N$ and, with the complete NS67 remainder,
\[
 R(t)=u\log t+v+\epsilon(t),\quad |\epsilon(t)|\le B/t,
 \quad q=u\log T+v,\quad B=\tfrac16\sum n|c_n|.
\]
For the chosen coefficients $d$, let
\[
 u_D=\alpha-\tfrac12\sum d_n,\qquad B_P=\tfrac1{12}\sum n|d_n|.
\]
Put $B_2(x)=x^2-x+1/6$. For $t\ge1$ the continuous periodic function
\[
 P(t)=-\tfrac12\sum n d_n B_2(\{t/n\})
\]
satisfies $P'=J-u_D$ almost everywhere and $|P|\le B_P$.
No period length is used in the bound. Write $f=Q/t$ and
$C_k=\int_{2^kT}^{2^{k+1}T}J(t)f(t)\,dt$. Integration by parts and the
triangle inequality, \emph{after} grouping, give
\[
 \sum_{k\ge0}|C_k|\le |u_D|\int_T^\infty|f|
 +B_P\left(|f(T)|+2\sum_{k\ge1}|f(2^kT)|+\int_T^\infty|f'|\right).
\]
The coefficient two counts shared endpoints. It may not be dropped.
Integrating the residual remainder first yields
\[
 Q(t)=\frac{u\log t+u+v}{t}+e(t),\qquad |e(t)|\le B/(2t^2).
\]
Moreover $f'=-R/t^3-Q/t^2$, so
\[
 |f'(t)|\le\frac{|2q+u|+2|u|\log(t/T)}{t^3}+\frac{3B}{2t^4}.
\]
Using the geometric sums $\sum_{k\ge1}4^{-k}=1/3$,
$\sum_{k\ge1}k4^{-k}=4/9$, and $\sum_{k\ge1}8^{-k}=1/7$, we obtain
\begin{align}
 \sum_{k\ge0}|C_k|\le\mathcal T_N(T):={}&
 |u_D|\left(\frac{|q+u|+|u|}{T}+\frac{B}{4T^2}\right)\\
 &+B_P\left[
 \frac{\tfrac53|q+u|+\tfrac89|u|\log2+
                    \tfrac12(|2q+u|+|u|)}{T^2}
 +\frac{8B}{7T^3}\right].
 \label{ns91-tail}
\end{align}
This bounds the sum of the absolute values of \emph{all} omitted dyadic
integrals, not just the absolute value of their signed sum. All constants
are explicit in the finite coefficients.

For $N=2^j$ and $T=2^L>N$, define, using the corrected representative,
\[
 \mathcal B_N(T)=
 \sum_{k=j}^{L-1}\left|\int_{2^k}^{2^{k+1}}
                  \widetilde J(t)Q(t)\frac{dt}{t}\right|
 +\mathcal T_N(T),\qquad b_N(T)=\alpha E-\mathcal B_N(T).
\]
If $b_N(T)>0$, then
\[
 \ell\ge b_N(T),\qquad E_N-E_{2N}\ge b_N(T)^2/K.
\]
The second conclusion follows from exact line search in the unchanged
projected direction. This is a finite sufficient inequality, not a new
RH criterion. No eventual bound on $\mathcal B_N(T)/E$ or $K$ is proved.

\paragraph{The regrouped absolute tail is also negligible cofinally.}
This statement concerns the tail, not the missing arithmetic gain.
The elementary identity
$\sum_{n\le x}\mu(n)\lfloor x/n\rfloor=1$ gives
$|\sum_{n\le x}\mu(n)/n|\le1+1/x$.
Define $S(x)=\sum_{n\le x}\mu(n)n^{-1}\log(x/n)$. Then
$s=S(2N)-S(N)$ and $S'(x)=x^{-1}\sum_{n\le x}\mu(n)/n$ almost everywhere, so
$|s|\le\log2+1/(2N)$. Consequently the endpoint-corrected coefficients obey
\[
 \sum|\widetilde d_n|\le3N\log2+\tfrac12,\qquad
 \sum n|\widetilde d_n|\le3N^2\log2+(\log2+\tfrac12)N.
\]
NS81 supplies $\|c\|_2^2\le32C_N$,
$C_N=48H_NN^3(N+2)$. Thus
\[
 |u|=O(N^{5/2}\sqrt{\log N}),\quad
 |v|=O(N^{5/2}\log^{3/2}N),\quad
 B=O(N^{7/2}\sqrt{\log N}),\quad
 |u_D|=O(N),\quad B_P=O(N^2).
\]
Insert these into \eqref{ns91-tail} at $T=N^6$:
\[
 \mathcal T_N(N^6)=O(N^{-5/2}\log^{3/2}N).
\]
The already proved NS83 bound $E_N\ge c/\log N$ for all sufficiently
large $N$, with no effective onset, then yields
\[
 \mathcal T_N(N^6)/E_N=O(N^{-5/2}\log^{5/2}N)=o(1/\log N).
\]
Therefore even the sum of absolute values of the \emph{grouped} outer
intervals can be controlled at the cofinal scale. This contrasts two
different estimates: it does not invalidate NS90's divergence of the
ungrouped divisor allowance at fixed $N$. The finite grouped arithmetic
integrals and true projected cost still require an independent estimate.

\paragraph{Exact endpoint evaluation.}
To avoid amplifying a loose endpoint remainder through many intervals,
the computation uses a closed expression for $Q(T)$. For $z>0$ and
$m=\lfloor z\rfloor$, direct cell integration gives
\[
 A(z)=z-m\log z+\log\Gamma(m+1),\qquad
 C(z):=\int_z^\infty\frac{\{x\}}{x^2}\,dx
 =H_m-\gamma-\log z+\frac{z-m}{z}.
\]
Indeed at an integer $m$, the complete cell sum telescopes to
$H_m-\gamma-\log m$; the partial cell gives the remaining term.
Integration by parts then shows
\[
 \int_T^\infty a_n(t)\frac{dt}{t^2}=\frac{A(T/n)}T+\frac{C(T/n)}n,
\quad Q(T)=\frac{\log T+1}T-
 \sum c_n\left(\frac{A(T/n)}T+\frac{C(T/n)}n\right).
\]
The Arb implementation uses $H_m-\gamma=\psi(m+1)$. No midpoint inverse,
uncertified quadrature, or omitted infinite tail enters the calculation.

\paragraph{Finite result and its limit.}
At $N=16,64,256$, the corrected dyadic bounds satisfy respectively
$b_N(T)/E>0.1793,0.2634,0.1867$ at $T=262144$.
At $N=256$ they guarantee relative gain $>0.02821$; the doubled-cutoff
replay gives $>0.02822$. These validate the new \emph{bound} on an already
known finite direction. They are not new best finite errors or a new
discovery of that direction's gain. In particular no gain lower bound is
asserted uniformly over $N$; the fractions above must not be extrapolated.

The arithmetic problem is now to estimate the grouped integrals uniformly,
and simultaneously control the true projected cost. Positivity of the
displayed finite dyadic terms is not assumed at other sizes; at $N=16$
several terms are negative. The cofinal estimate needed for RH remains
open. The first-primitive identity and endpoint correction alone do not
supply it.

\paragraph{Uniform-estimate attempt: what the quadratic rewrite actually says.}
The divisor inversion is already NS78's identity, not a new construction.
For arbitrary real $\beta,c_1,\ldots,c_N$, extend $c$ by zero and put
\[
 R_{\beta,c}=\beta\tau-\sum_{n\le N}c_na_n,\quad
 h_k=\langle R_{\beta,c},a_k\rangle,\quad
 \delta_m=\beta\mathbf1_{m=1}+\sum_{n\mid m}c_n.
\]
Dirichlet convolution gives $\mu*\delta=\beta\mu+c$.
With $w_k=\log(2N/k)$, define
\[
 \mathfrak q_N(R_{\beta,c})
 =-\sum_{dm\le2N}\mu(d)w_{dm}\delta_m h_{dm},\qquad
 \ell_N(R_{\beta,c})=-\sum_{N<k\le2N}\mu(k)w_k h_k.
\]
These are respectively a quadratic and a linear expression in
$(\beta,c)$. Direct substitution gives the exact identity
\[
 \mathfrak q_N(R_{\beta,c})=\beta\ell_N(R_{\beta,c})
 -\sum_{k\le N}w_k h_k(\beta\mu(k)+c_k).
\]
Thus the old normal equations $h_k=0$ for $k\le N$ do remove the final
sum. But they also force $c=\beta G_N^{-1}g_N$, where
$g_N=(\langle\tau,a_k\rangle)_{k\le N}$, and
$R_{\beta,c}=\beta R_N$, because the old Gram is positive definite.
Consequently, on this constrained space,
\[
 \mathfrak q_N(R_{\beta,c})=\beta^2\ell_N(R_N),\qquad
 \|R_{\beta,c}\|^2=\beta^2 E_N.
\]
A uniform lower bound $\mathfrak q_N\ge\varepsilon_N\|R_{\beta,c}\|^2$
on these residuals is therefore exactly the missing scalar bound
$\ell_N(R_N)\ge\varepsilon_NE_N$. Writing it as a quadratic expression
does not establish its sign. No new positivity or impossibility theorem
is inferred; estimating its actual arithmetic terms remains possible.
The finite Maxima checks verify inversion through index eight and the
old-equation reduction only, not an eventual inequality.

\paragraph{Precise next estimate.}
For the corrected representative define
$\rho_N=\mathcal B_N(N^6)/E_N$ and $\kappa_N=K_N/E_N$.
This selected-direction approach would succeed if, for $N=2^j$ eventually,
\[
 \rho_N<\log2,\qquad
 \frac{(\log2-\rho_N)^2}{\kappa_N}\ge\frac a j
 \quad(a>0).
\]
Both the finite interior margin and its relation to the true projected
cost need proof; neither has been bounded uniformly here. The stronger
separate hypotheses $\rho_N\le\log2-\varepsilon$ and $K_N\le K_0$,
with fixed positive $\varepsilon,K_0$, would suffice by NS83, but are not
assumed. The full NS87 observation-energy route is still available if
this single predetermined direction does not admit the needed estimate.
